% !TeX root = ../main.tex

\chapter{Preparation}

\section{Outline}

\begin{description}
  \item [Free electron gas]
  Drude model; Sommerfeld model; magneto-transport
  \item [Electrons in a crystal]
  Crystal structure (such as symmetry); Diffraction experiments; Lattice vibrations
  \item [Band theory]
  Bloch's theorem; Nearly free electron model (simplified model); Tight binding Model (simplified model); Applications of band theory; Semiconductor devices
  \item [Exotic band theory]
  Dirac band; [Topological band; Flat band]
  \item [Beyond band theory (Many-body physics)]
  Magnetism; Superconductivity
\end{description}

\section{Course Overview}

\begin{description}
  \item [Lecture 1] Introduction
  \item [Lecture 2] Free electron model of metals
  \item [Lecture 3] Lattice vibrations and phonons
  \item [Lecture 4] Crystal symmetry and crystal structure
  \item [Lecture 5] Scattering experiments (X-ray diffraction and electron diffraction)
  \item [Lecture 6] Electron waves in crystals (Bloch's theorem)
  
  \emph{The two models help us to Understand why there's a gap in some crystal}
  \item [Lecture 7] Band theory I: Nearly free electron model
  \item [Lecture 8] Band theory II: Tight binding model
  \item [Lecture 9] Band theory III: Fermi surface ad effective mass
  \item [Lecture 10] Physics of semiconductors
  \item [Lecture 11] Electrons in a magnetic field: Quantum Hall effect
  \item [Lecture 12] Dirac band theory I: band structure of graphene
  \item [Lecture 13] Dirac band theory II: Landau quantization of graphene
  \item [Lecture 14] Magnetism
  \item [Lecture 15] Superconductivity
  \item [Lecture 16] Topological Physics
\end{description}

\section{References}

\begin{itemize}
  \item Introduction to Solid State Physics, Charless Kittel
  \item Solid State Physics, Neil Ashcroft and David Mermin
  \item The Oxford Solid State Basics, Steven H. Simon
  \item Solid State Physics: Essential Concepts, David W. Snoke
\end{itemize}